\documentclass[10pt,letterpaper]{article}
\usepackage[margin=0.45in]{geometry}
\usepackage{amsmath,amssymb,microtype,enumitem}
\usepackage[T1]{fontenc}
\pagestyle{empty}
\setlength{\parindent}{0pt}
\setlength{\parskip}{2.5pt}
\setlength{\abovedisplayskip}{4pt}
\setlength{\belowdisplayskip}{4pt}
\setlist{nosep,leftmargin=10pt}
\begin{document}
\begin{center}
{\large\bfseries Prediction-Aware Robust Convex Optimization for\\ Energy-Harvesting Wireless Packet Scheduling}\\[3pt]
{\bfseries Arghadeep Saha \quad and \quad Protyay Dey }\\[2pt]
{\small EE 541: Convex Optimization for Communications, Networking, and Machine Learning\\ University at Buffalo, SUNY \quad $\vert$ \quad Fall 2026 Project Proposal}
\end{center}
\vspace{-7pt}
\textbf{Research question and motivation.} How can finite-horizon convex optimization be used to manage the scheduling of packets and the expenditure of power in an energy harvesting wireless sensor in light of uncertainties about future channel states, packet arrivals, and energy harvests? The sensor has a finite battery and buffer, has channels that vary with time, and harvests energy intermittently. Sending a packet immediately clears the queue but uses up energy which would have been better utilized in future when channel conditions become favorable. Not sending now means that there is a risk of buffer overload and packet loss.

\textbf{System model and formulation.} For slots $t=0,\ldots,T-1$, the decisions are transmission power $p_t$, transmitted data $d_t$ (continuous packet equivalents), queue occupancy $b_t$, battery energy $E_t$, rejected arrivals $o_t$, and spilled energy $w_t$. Inputs are channel power gain $h_t\geq0$, arrivals $l_t\geq0$, and harvested energy $H_t\geq0$. Given initial states $b_0,E_0$ and fixed physical parameters, solve

For slots $t=0,\ldots,T-1$, the actions consist of transmission power $p_t$, transmitted data $d_t$ (continual packet equivalents), number in queue $b_t$, battery energy $E_t$, rejected arrivals $o_t$, and wasted energy $w_t$. The inputs include the channel power gain $h_t\geq0$, arrivals $l_t\geq0$, and harvested energy $H_t\geq0$. With given initial states $b_0,E_0$ and given physical constants, find

\begin{equation}
\begin{aligned}
\min_{p,d,b,E,o,w}\quad &\sum_{t=0}^{T-1}\left(b_t+\eta o_t+\lambda\tau p_t\right)+\beta b_T\\[-1pt]
\text{s.t.}\quad &d_t\leq\frac{B\tau}{L}\log_2\!\left(1+\frac{h_tp_t}{N_0B}\right),\\[-1pt]
&b_{t+1}=b_t-d_t+l_t-o_t,\qquad E_{t+1}=E_t-\tau p_t+H_t-w_t,\\[-1pt]
&0\leq d_t\leq b_t,\quad d_t\leq D_{\max},\quad 0\leq o_t\leq l_t,\quad w_t\geq0,\\[-1pt]
&0\leq b_t\leq N_b,\quad 0\leq E_t\leq E_{\max},\quad 0\leq p_t\leq P_{\max},\quad \tau p_t\leq E_t,\\[-1pt]
&0\leq b_T\leq N_b,\qquad 0\leq E_T\leq E_{\max}.
\end{aligned}
\end{equation}
Here $B$ is bandwidth, $N_0$ noise spectral density, $L$ packet size, and $\tau$ active transmission time; $\eta,\lambda,\beta\geq0$ are cost weights. The scheduling interval is $\Delta=5$ s, with active transmission time $\tau\leq\Delta$. Data arrivals are admitted after service; rejected arrivals obey $0\leq o_t\leq l_t$. Harvested energy is usable only from the next interval, and overflow of the energy store is captured by $w_t$. The rate function is concave in $p_t$ for fixed $h_t\geq0$; its hypograph and the remaining affine constraints define a convex feasible set. The objective is convex for nonnegative penalty weights, permitting globally optimal finite-horizon solutions.

\textbf{Method and EE 541 extension.} Receding-horizon MPC will be implemented in Python (NumPy and CVXPY) with the following approach: at each slot, observe the actual state of queue/battery/channel, solve the problem with the forecasting horizon, do the first action and update based on realization. The planning horizon will be taken equal to 10-20 slots, only the first action will be taken from the result of optimization, and then the optimization will be run again using the realization of the state. For the EE 541 extension, we will compare the nominal MPC and conservative scheduler with the substitution of the predicted gain $\widehat h_t$ with $\underline h_t=\max\{0,\widehat h_t-\delta_t\}$ in the rate constraint. We will consider conservative bounds on the harvested energy and arrivals. With fixed uncertainty bounds, convexity is preserved.

\textbf{Evaluation and verification.}
We will consider a single energy harvesting sensor with buffer size $N_b=32$, maximum number of packet transmissions $D_{\max}=10$, and random packet arrivals $l_t\in\{0,\ldots,10\}$ per time unit. We will analyze performance under varying channel states, energy harvesting rates, and forecast uncertainties. We will compare fixed power control, greedy policy, channel-aware policy, nominal MPC, robust MPC, adaptive MPC, and adaptive robust MPC against an off-line optimal scheme with the same simulations traces. We will consider average queueing delay, battery usage, buffer overflows, throughput, energy usage, and computational time. In addition, we will analyze adaptive penalty dynamics, trade-off between delay and energy consumption, and forecast error sensitivity. Moreover, we will check compliance with CVXPY convexity requirements, solver convergence, and feasibility of the constraints.

\textbf{Expected outcome.} An implemented simulator and CVXPY program, along with tested constraints and data related to the tradeoff between prediction accuracy, queueing delays, dropped packets, and energy consumed in transmission. It will not be assumed that robust scheduling works well for all scenarios.

\textbf{Team contributions.} \textbf{Protyay Dey:}  uncertainty-aware extension, mathematical verification and analysis. \textbf{Arghadeep Saha:} convex formulation, wireless/harvesting simulator, baseline implementations, and experiment visualization. Both members will validate results and prepare the report and presentation.

\textbf{References.} \textbf{[1]} O. Ozel et al., ``Transmission with Energy Harvesting Nodes in Fading Wireless Channels: Optimal Policies,'' \emph{IEEE JSAC}, 2011. \textbf{[2]} J. Yang and S. Ulukus, ``Optimal Packet Scheduling in an Energy Harvesting Communication System,'' \emph{IEEE Trans. Commun.}, 2012. \textbf{[3]} S. Boyd and L. Vandenberghe, \emph{Convex Optimization}, 2004. \textbf{[4]} M. Shi, X. Lin, and S. Fahmy, ``Competitive Online Convex Optimization with Switching Costs and Ramp Constraints,'' \emph{IEEE/ACM Trans. Netw.}, 2021.
\end{document}